\section{From X25519 to a secure protocol}

X25519 is neither encryption nor authentication. It computes raw shared key
material. An active attacker can substitute public keys and complete two
independent exchanges:

\begin{center}
\begin{tikzpicture}[
  node distance=27mm,
  participant/.style={draw=black!25,rounded corners,fill=black!2,minimum width=18mm,minimum height=8mm},
  >={Stealth[length=2mm]}
]
\node[participant] (alice) {Alice};
\node[participant,right=of alice,draw=runningorange] (mallory) {Mallory};
\node[participant,right=of mallory] (bob) {Bob};
\draw[<->] (alice) -- node[above]{secret $S_{AM}$} (mallory);
\draw[<->] (mallory) -- node[above]{secret $S_{MB}$} (bob);
\node[below=8mm of mallory,text=runningorange] {$S_{AM}\ne S_{MB}$ in general};
\end{tikzpicture}
\end{center}

Secrecy against a passive observer is not proof of the peer's identity.
Authentication must come from a certificate and signature, a pre-shared key,
an authenticated static key, or another mechanism defined by the protocol.

\subsection{Key derivation and context binding}

A raw Diffie--Hellman coordinate is not an application key. A construction
such as HKDF first extracts a pseudorandom key and then expands labeled keys
for specific purposes \cite{rfc5869}:
\[
  \text{raw DH}
  \longrightarrow
  \text{all-zero handling}
  \longrightarrow
  \operatorname{HKDF}(S,A,B,\text{roles},\text{transcript})
  \longrightarrow
  \begin{cases}
  K_{A\to B},\\
  K_{B\to A}.
  \end{cases}
\]
Length prefixes or another unambiguous encoding prevent distinct transcripts
from collapsing to the same KDF input. Ordered public keys and role labels stop
reflection and unknown-key-share confusions. Separate directional keys simplify
nonce and replay reasoning.

RFC~7748 permits a generic implementation to check for all-zero output, while
newer profiles can strengthen that rule. The X25519-based HPKE suite requires
the recipient to abort on an all-zero DH result and uses labeled HKDF inputs
\cite{rfc9180}. A high-level library should surface an unacceptable-peer-key
error rather than silently continuing.

After key derivation, an AEAD supplies confidentiality and ciphertext
integrity. Its nonce requirements remain independent protocol obligations.

\subsection{Forward secrecy, signatures, and quantum attacks}

Fresh ephemeral X25519 keys can provide forward secrecy when the handshake is
authenticated and the ephemeral secrets are erased. Curve25519 itself does not
guarantee freshness or erasure.

Ed25519 is a signature scheme over a related Edwards representation
\cite{rfc8032}. It is not X25519. Its public-key encoding, private-key
derivation, and security purpose differ. Conversion routines exist in some
libraries, but distinct purpose-specific key pairs are the safer default.

Finally, Curve25519 is not post-quantum. A sufficiently capable quantum
computer running Shor's algorithm would solve ECDLP \cite{shor1997}. Hybrid
deployments combine classical authentication and X25519 with a standardized
post-quantum mechanism according to a reviewed protocol; ad hoc concatenation
is not a protocol.

\subsection{Composition checklist}

For production work:
\begin{enumerate}
  \item use a maintained X25519 implementation rather than handwritten field
        arithmetic;
  \item follow the enclosing protocol's all-zero rule;
  \item bind ordered public keys, roles, suite identifiers, and the transcript
        into a KDF;
  \item authenticate the peer;
  \item derive separate keys by purpose and direction;
  \item use an AEAD and obey its nonce limits;
  \item prefer a reviewed construction such as TLS, Noise, or HPKE.
\end{enumerate}

\begin{takeaway}
The chain is modular: efficient scalar multiplication supplies the public
operation; a CDH-style assumption says an outsider cannot combine two public
points into their shared point; solving ECDLP is one sufficient break. X25519
packages this group action, a KDF turns its output into keys, authentication
binds those keys to identities, and an AEAD protects messages.
\end{takeaway}
